\documentclass{uflamon}          % classe base para a monografia

%==============================================================================
% Utilizacao de pacotes
\usepackage[T1]{fontenc}         % usa fontes postscript com acentos
\usepackage[brazil]{babel}       % hifenização e títulos em português do Brasil
\usepackage[utf8]{inputenc}     % permite edição direta com acentos
\usepackage{amsmath}             % pacote da AMS para Matemática Avançada
\usepackage{amssymb}             % símbolos extras da AMS
\usepackage{latexsym}            % símbolos extras do LaTeX
\usepackage{graphicx}            % para inserção de gráficos
\usepackage{listings}            % para inserção de código
\usepackage{fancyvrb}            % para inserção de saídas de comandos
%\usepackage{enumerate}           % para personalizar lista enumeradas 
											%(incluso na classe)
\usepackage{longtable}           % para tambelas muito grandes NOVO!!!!

\usepackage{colortbl} % cores em tabelas
\newcolumntype{Z}{|>{\columncolor[gray]{0.9}}l|} %cor cinza em células
%\usepackage{array} % já incluso na classe
\newcolumntype{L}[1]{>{\raggedright\let\newline\\\arraybackslash\hspace{0pt}}m{#1}}
\newcolumntype{C}[1]{>{\centering\let\newline\\\arraybackslash\hspace{0pt}}m{#1}}
\newcolumntype{R}[1]{>{\raggedleft\let\newline\\\arraybackslash\hspace{0pt}}m{#1}}
\usepackage{multirow} % para juntar duas linhas em uma só
\usepackage{multicol} % para uso de várias colunas
%\usepackage{subfig}

% cores para os links cruzados
\usepackage{color}
\definecolor{rltred}{rgb}{0.2,0,0}
\definecolor{rltgreen}{rgb}{0,0.2,0}
\definecolor{rltblue}{rgb}{0,0,0.2}

\usepackage[colorlinks=true,
            urlcolor=rltblue,       % \href{...}{...} external (URL)
            filecolor=rltgreen,     % \href{...} local file
            linkcolor=black,       % \ref{...} and \pageref{...}
            citecolor=black,
            pdftitle={Determinação da Perda de Carga por Acessórios},
          pdfauthor={Marcus Bruno Fernandes Silva},
          pdfsubject={Determinação da Perda de Carga por Acessórios},
          pdfkeywords={perda de carga}%
]{hyperref} % para referência cruzadas
%\usepackage{hyperref}            % para referência cruzadas
\usepackage{subfigure}           % figuras dentro de figuras
\usepackage{caption}            % remodelando o formato dos títulos de 
                                 % tabelas e figuras

% configuração padrão do listings   
\lstset{%
   language=Java,
   extendedchars=true,
   tabsize=3,
   basicstyle=\footnotesize\ttfamily,
   stringstyle=\em,
   showstringspaces=false 
}

%abnt-emphasize=bf coloca o título das bibliografias em negrito
%abnt-thesis-year=both
\usepackage[alf,num,abnt-etal-cite=3,abnt-etal-list=3,abnt-url-package=url,abnt-emphasize=bf]{abntex2cite}

% redefinindo formatação de títulos de tabelas e figuras

%==============================================================================
\usepackage{microtype} 			% para melhorias de justificação
\usepackage{float}
\usepackage{titling}
\usepackage[version=4]{mhchem}
\usepackage{array}
\usepackage{booktabs}
%\usepackage[authoryear]{natbib}

% Especificando hifenizações que por ventura LaTeX não saiba fazer
% Por padrão 99,9% dos termos em português devem ser hifenizados corretamente.
\hyphenation{hardware software Li-nux am-bien-te diag-nos-ti-car coor-de-na-ção 
FAE-PE Recovery TelEduc Williams UFLA Sérgio Renata Reynolds disciplina}

%==============================================================================
% Dados da monografia, capa: autor, titulo, banca, etc... - SUBSTITUA DE ACORDO
%==============================================================================
\author{Felipe Alexander Julio \\
       Gabriela Nunes \\
       Marcus Bruno Fernandes Silva \\
       Pétrus Damasceno}
\subtitle{Exemplo para os Usuários}
%\author{Joao}
\title{Determinação da Perda de Carga por Acessórios}
\engtitle{Use of Uflamon Class}
%\engsubtitle{Sample for Users}
%\edicao{3$^a$ edição revista, atualizada e ampliada}
\date{2017}
\tipo{Pré-projeto entregue e apresentado à Professora Renata de A. B. L. Corrêa como requisito da disciplina GNE334.}
\orientador{Professora Renata de A. B. L. Corrêa}
\newcommand{\orient}{Professora Renata de A. B. L. Corrêa}
\newcommand{\otipo}{Pré-projeto entregue e apresentado à Professora Renata de A. B. L. Corrêa como requisito da disciplina GNE334.}
\newcommand{\olocal}{Lavras -- MG}
\local{Lavras -- MG}

%##################################################

\newcommand \fazercapa{
\begin{titlepage}
    \centering

    \includegraphics[scale=0.45]{logoufla.jpg}%

    \vspace{1cm}

    \Large\theauthor
    
    \vspace{\stretch{1}}

    \bfseries\LARGE\MakeUppercase\thetitle

    \vspace{\stretch{1}}

    \Large\MakeUppercase\olocal\\
    \MakeUppercase\thedate
    \null
    \cleardoublepage
\end{titlepage}
}

\newcommand \fazerfolhaderosto{
    \begin{titlepage}
        \newpage

        \begin{center}
            \MakeUppercase\theauthor

            \vspace{\stretch{1}}

            \textbf{\MakeUppercase\thetitle}

            \vspace{\stretch{1}}
            
            \hfill\begin{minipage}{8cm}{\otipo}\end{minipage} % tipo da monografia

            \vspace{\stretch{.58}}

            \ifthenelse{\equal{\@orientador}{\@empty}} % % se orientador nao existir,
            {}%                                            % nao produza nada
            {\orient\\
             Orientadora
            }

            \vspace{\stretch{.86}} 

            \textbf{\MakeUppercase\olocal}\\ 
            \textbf{\MakeUppercase\thedate}
            \end{center}
    \end{titlepage}
}

\newcolumntype{P}[1]{>{\centering\arraybackslash}p{#1}}
\graphicspath{{./figuras/}}

\newcommand{\tabitem}{~~\llap{\textbullet}~~}
\newcommand{\R}{$^{\tiny{\textregistered}}$~}
%==============================================================================
% Começar o documento

%\newsubfloat{figure}

%\includeonly{capitulos/materiais_metodos,capitulos/resultados}
\newcommand{\pa}[1]{\left(#1 \right)}
\begin{document}

\fazercapa

\fazerfolhaderosto

\tableofcontents

\clearpage

\pagestyle{ufla}

%==============================================================================
% incluindo os capitulos

\newcommand{\citeay}[1]{\citeauthoronline{#1} (\citeyear{#1})}
\newcommand{\citeayp}[1]{(\citeauthoronline{#1} \citeyear{#1})}
\newcommand{\citean}[1]{\citeauthoronline{#1} \cite{#1}}

\chapter{Introdução}\label{intro}

\section{Referencial Teórico}


Segundo \citeay{Munson04} ``a mecânica dos fluidos é a parte 
da mecânica aplicada que se dedica à
análise do comportamento dos líquidos e gases, tanto em equilíbrio quanto em 
movimento. Obviamente, o escopo da mecânica
dos fluidos abrange um vasto conjunto de 
problemas.''

Na análise do escoamento de fluidos é necessário o 
balanço de massa e de momento para desenvolvimento de equações
apropriadas aos fenômenos a serem analisados \citeayp{Cremasco14}.
A descrição matemática da fluidodinâmica é dada pelas 
equações de massa \eqref{eqcontinuidade}
e equações do movimento \eqref{eqQDM}.


\begin{equation}\label{eqcontinuidade}
    \frac{\partial\rho}{\partial t} + 
    \vec{\nabla} \cdot \rho \textbf{u}  = 0
\end{equation}

\begin{equation}\label{eqQDM}
    \rho \left(
        \frac{\partial \rho}{\partial t} + 
    \textbf{u} \cdot \vec{\nabla}\rho\textbf{u}
    \right)
    = 
    \rho\vec{\nabla}\varphi -
    \vec{\nabla} P +
    \vec{\nabla}\cdot \tau
\end{equation}

onde $\rho$ é a massa específica do fluido, \textbf{u} é a sua velocidade,
$\vec{\nabla}\varphi$ é o vetor intensidade do campo, $\vec{\nabla} P$ é o gradiente
de pressão e $\tau$ é a tensão extra, dinâmica ou viscosa.

Ao manipular as equações \eqref{eqcontinuidade} e 
\eqref{eqQDM}, considerando que o fluido se encontra em
regime permanente, é invíscido e incompressível, chega-se a
equação simplificada para a energia mecânica:

\begin{equation}\label{eqBern}
    \left(
        z + \frac{P}{\rho g} + \frac{u^2}{2g}
    \right)_{1}
    =
    \left(
        z + \frac{P}{\rho g} + \frac{u^2}{2g}
    \right)_{2} + h_L

\end{equation}

onde os subíndices $1$ e $2$ indicam dois pontos de escoamento dinstintos, e $h_L$ é
a perda de carga dada em unidades de comprimento.

Conforme \citeay{Cremasco14}, um fluido em um sistema de 
escoamento que apresente 
acessórios estará sujeito a uma perda de
carga, uma vez que, segundo o autor, quando o fluido entra em contato com 
um acessório ocorre a separação de uma camada
do escoamento, formando correntes turbulentas que dissipam energia mecânica. 
Essa energia dissipada representa a conversão irreversível de energia mecânica 
em energia térmica não desejada e a perda de energia por transferência de 
calor \citeayp{Fox14}.


A perda de carga oriunda de um acessório em um sistema de escoamento é definida por meio da equação:

\begin{equation}
    h_L = k_f \frac{u^2}{2g}
\end{equation}

onde $k_f$ é o coeficiente de perda de carga localizada característico do acessório, 
$u$ é a velocidade média do fluido e $g$ é a aceleração local da gravidade. 

Ainda conforme o autor, $k_f$ é constante apenas em regime turbulento. Podem ser 
encontradas na literatura tabelas que apresentam os valores de $k_f$ para válvulas
e acessórios específicos.

De acordo com \citeay{Potter15}, o regime turbulento é caracterizado 
pelo movimento
aleatório das partículas do fluido na
trajetória do escoamento. Tal regime pode ser previsto por um número 
adimensional denominado por número de Reynolds,
o qual é calculado, para tubulações cilíndricas, conforme:

\begin{equation}
    Re = \frac{u D}{\nu}
\end{equation}
 
onde $D$ é o diâmetro característico da tubulação e $\nu$ é a viscosidade 
cinemática do fluido.

Ainda conforme \citeay{Potter15}, valores para o número de 
Reynolds acima de 2000 são 
características de um escoamento em regime
turbulento, sendo este o valor mínimo usado na maioria das aplicações em engenharia.
Entretanto, tal valor pode variar em decorrência da geometria e da 
rugosidade das paredes da tubulação.

Portanto, além do cálculo adimensional, pode ser realizada uma análise 
qualitativa por meio do despejo de partículas
corantes na corrente de escoamento, avaliando se ocorre uma mistura imediata, ou 
não. Caso não, o regime é laminar.

\section{Aplicabilidade}

O cálculo da perda de carga é necessário em qualquer tipo de planejamento
e instalação de tubulações, uma vez que o mesmo está relacionado diretamente
com o matérial utilizado e as características da tubulação.

 Algumas situações que requerem uma análise da perda de carga podem ser 
exemplificadas como: possível cavitação em
bombas, danos à tubulação de edifícios, insuficiência na captação de água e/ou 
esgoto para tratamento e controle de
fluxo em reatores e outros 
equipamentos da indústria química.

\section{Objetivos}

O objetivo desse projeto é determinar a perda de carga e o coeficiente de 
perda de carga localizada em um tubo conectado
a dois reservatórios com um fluxo de água constante através de acessórios 
(individualmente) como registros e cotovelos a
partir da variação de pressão entre os pontos anterior e 
posterior aos acessórios.

\chapter{Motivação}\label{motivacao}
Para o desenvolvimento de um projeto industrial que envolva transporte de 
fluidos é necessário o conhecimento do
estado energético-mecânico do fluido para que o mesmo possa ser deslocado 
conforme a necessidade do sistema.  Avaliando,
neste caso, a necessidade ou não da utilização de sistemas fluidomecânicos 
(bombas, compressores, por exemplo). 

A aferição do estado energético-mecânico do fluido em um determinado ponto de 
uma tubulação, por exemplo, requer a
investigação precisa dos pontos na tubulação por onde passou e o estado 
inicial do mesmo. Desta forma, elencar-se-á as
perdas de carga envolvidas durante o processo e, poder-se-á, então, determinar 
o estado de um fluido em um ponto e
avaliar se será necessário a injeção, ou remoção, de carga (energia) ao 
sistema. A presença de acessórios em uma
tubulação é quase que inevitável, seja por necessidades de otimização de 
espaço, seja por segurança ou por qualquer
outro motivo que seja necessária a instalação de uma válvula, por exemplo. 
Portanto, a investigação da perda de carga em
um escoamento é vital ao desenvolvimento de um projeto que atenda, com 
qualidade e segurança, o objetivo do sistema.

%%% Local Variables:
%%% mode: latex
%%% TeX-master: "../main_archive"
%%% End:

\chapter{Descrição do Projeto}\label{descricao}

\section{Materiais Utilizados}

 \begin{table}[H]\label{t:custos}
     \centering
     \begin{tabular}{l c }
         \toprule
         \textbf{Material} & \textbf{Valor (R\$)} \\
         \midrule
         Adaptador (conexão entre balde e cano) & 1,00 - 5,00 \\
         Baldes & 20,00 - 30,00 \\
         Canos PVC & 25,00 - 50,00 \\
         Cola Durepoxi & 10,00 - 20,00 \\
         Cotovelos 90$^\circ$ & 1,00 - 5,00\\
         Cronômetro & 10,00 - 20,00 \\
         Registro rosqueável & 30,00 - 50,00 \\
         Suportes para a estrutura (madeira) & 20,00 - 50,00 \\
         Trena & 10,00 - 15,00 \\
         Tubos transparentes (manômetro) & 15,00 - 20,00 \\
         Válvula esférica & 30,00 - 50,00 \\
         \midrule
         \textbf{TOTAL} & 172,00 - 315,00\\
         \bottomrule
     \end{tabular}
     \caption{Custos}
 \end{table}

\section{Montagem}

A estrutura se inicia com um recipiente posicionado a uma elevação 
considerável do nível convencionado como $z = 0$. Esse
recipiente será um balde de alta capacidade volumétrica, o qual possuirá um ``ladrão''
na sua cota máxima, afim de evitar-se transbordamento do fluido excedente e também
estará ligado, em sua parte inferior, 
à tubulação de PVC através de um
adaptador, e após essa conexão será instalado um registro para controle da 
saída do fluido. A seguir, a tubulação fará
uma curva de $90^\circ$, descendo até o nível $z = 0$, 
onde fará uma nova curva de $90^\circ$
e prosseguirá reta até o local onde será
acoplado o acessório escolhido. Inicialmente, será analisada a perda de carga 
em uma válvula esférica e para tal é
necessário medir a diferença de pressão no tubo antes e depois da
válvula, com o auxílio de um manômetro.

Manômetros são instrumentos utilizados para determinação da diferença de 
pressão entre dois pontos de um fluido através
da diferença de elevação entre estes dois pontos. Dessa forma, o manômetro 
será constituído de dois tubos transparentes
acoplados ao tubo de PVC, um deles um pouco antes da válvula e o outro um 
pouco depois. Nesses tubos serão medidas as colunas de líquido. 

Em seguida, a estrutura irá terminar com a tubulação despejando o fluido em um 
outro recipiente de capacidade
volumétrica conhecida de forma que, através da cronometragem do tempo que esse 
recipiente leva para completar seu volume, seja possível 
determinar a vazão volumétrica. 




%==============================================================================

% Incluindo bibliografia
%\bibliographystyle{plain}             % estilo para labels em numeros
%\bibliographystyle{alpha}             % estilo para labels em iniciais
\bibliographystyle{abntex2-num}         %% estilo para referências usando ABNT, 
 %para usar o abntex2-alf precisa arrumar o nome dos autores na bibliografia
%\bibliographystyle{abntex2-alf}           % estilo para referências usando ABNT, 

%inclui Referências Bibliográficas
\referencias
\bibliography{bibliografia/b}

%==============================================================================
% Fim do texto
\end{document}
